%% maestro2tex -- generated figure; do not edit by hand.
\providecommand{\MaestroRoot}{include/analog/}
\begin{figure}[htbp]
  \centering
  {\pgfplotsset{compat=1.18}%
  \begin{tikzpicture}
    \begin{axis}[
      width=0.82\linewidth, height=6.2cm,
      xlabel={Frequency~[\si{\hertz}]}, ylabel={Loop gain $|T|$~[\si{\decibel}]},
      grid=both,
      major grid style={line width=0.2pt, draw=black!14},
      minor grid style={line width=0.2pt, draw=black!7},
      minor tick num=1,
      axis line style={line width=0.4pt, draw=black!45},
      tick style={line width=0.4pt, draw=black!45},
      tick align=outside,
      label style={font=\small}, tick label style={font=\small},
      enlarge x limits=0.02, enlarge y limits=0.10,
      every axis plot/.append style={line join=round, no marks},
      legend style={font=\footnotesize, draw=black!25, fill=white,
                    fill opacity=0.9, text opacity=1, row sep=1pt},
      legend cell align=left, legend pos=south west,
      xmode=log,
      scaled y ticks=false,
      y tick label style={/pgf/number format/fixed,
                          /pgf/number format/precision=0}
    ]
      \addplot[mtxA, line width=1.1pt, solid] table {\MaestroRoot data/px_ctrl_loopgain_00.dat};
      \addlegendentry{tt}
      \addplot[mtxB, line width=1.1pt, dashed] table {\MaestroRoot data/px_ctrl_loopgain_01.dat};
      \addlegendentry{ff}
      \addplot[mtxC, line width=1.1pt, dotted] table {\MaestroRoot data/px_ctrl_loopgain_02.dat};
      \addlegendentry{fs}
      \addplot[mtxD, line width=1.1pt, dashdotted] table {\MaestroRoot data/px_ctrl_loopgain_03.dat};
      \addlegendentry{sf}
      \addplot[mtxE, line width=1.1pt, densely dashed] table {\MaestroRoot data/px_ctrl_loopgain_04.dat};
      \addlegendentry{ss}
    \end{axis}
  \end{tikzpicture}}
  \caption[{Potentiostat control-loop gain over five process corners, measured through the \SI{0}{\volt} probe\,\dots}]{Potentiostat control-loop gain over five process corners, measured through the \SI{0}{\volt} probe \texttt{IPRB1} in the RE feedback path (Figure~\ref{fig:tb-pixel}). $R_f = \SI{35}{\kilo\ohm}$, cell $R_{\mathrm{ce}} = R_u = \SI{1}{\kilo\ohm}$, $R_{\mathrm{ct}} = \SI{10}{\kilo\ohm}$, $C_{\mathrm{dl}} = \SI{100}{\nano\farad}$, \SI{2.5}{\volt} supply. This is the full loop -- A1, A2, both feedback networks and the electrode together -- and is not the open-loop amplifier gain of Section~\ref{ss:ampab}. \texttt{tt} is at \SI{37}{\celsius}, the four skew corners at \SI{25}{\celsius}. The corner separation drawn here is a lower bound: the corner set holds the passive model rows at typical (pinned-passive caveat, page~\pageref{p:pinned-passives}).}
  \label{fig:px-ctrl-loopgain}
\end{figure}
